\chapter*{What a trace forgets}
\markboth{What a trace forgets}{What a trace forgets}
\phantomsection
\addcontentsline{toc}{chapter}{What a trace forgets}
\label{v4-ch:trace-prologue}

A trace assigns a number to an operation. To recover the operation from
that number, one must know what the trace has forgotten. Four elementary
calculations will guide the arithmetic constructions in this volume.
They retain, respectively, a primitive of a relation, a higher product,
an intermediate module state, and a geometric condition on an arithmetic
operator. Each calculation identifies data that a comparison must retain.

\section*{Gaussian reduction and multiplication}
\label{v4-prologue:gaussian}

Fix a real number \(\hbar>0\). The normalized Gaussian functional on
polynomials is
\[
 \mathbb E_\hbar[p]
 =\frac{1}{\sqrt{2\pi\hbar}}
   \int_{\mathbb R}p(x)e^{-x^2/(2\hbar)}\,dx.
\]
Integration by parts has no boundary term and gives
\begin{equation}\label{v4-prologue:eq:gaussian-relation}
 \mathbb E_\hbar[xp-\hbar p']=0.
\end{equation}
This relation defines a two-term cochain complex
\[
 C^{-1}=\mathbb R[x]\varepsilon,
 \qquad C^0=\mathbb R[x],
 \qquad d(p\varepsilon)=xp-\hbar p'.
\]
There are no other degrees. Its differential is injective: if \(p\)
has degree \(n\) and nonzero leading coefficient, then \(xp-\hbar p'\)
has degree \(n+1\) with the same leading coefficient. The relations
\[
 [x]=0,\qquad
 [x^{n+1}]=\hbar n[x^{n-1}],\qquad n\ge1,
\]
show that \(H^0(C)\) is spanned by \([1]\). Equation
\eqref{v4-prologue:eq:gaussian-relation} and
\(\mathbb E_\hbar[1]=1\) show that \([1]\ne0\), so
\[
 H^{-1}(C)=0,\qquad H^0(C)\cong\mathbb R.
\]
Under this identification,
\begin{equation}\label{v4-prologue:eq:moments}
 \mathbb E_\hbar[x^{2r+1}]=0,
 \qquad
 \mathbb E_\hbar[x^{2r}]
 =\hbar^r\frac{(2r)!}{2^r r!}.
\end{equation}
The recurrence proves the formula without evaluating another integral.

Polynomial multiplication does not descend to this quotient. Indeed,
\[
 d\varepsilon=x,\qquad
 d(x\varepsilon)=x^2-\hbar,
\]
so \([x]=0\) but \([x^2]=\hbar[1]\ne0\). Equivalently,
\(\mathbb E_\hbar[x]^2=0\) and
\(\mathbb E_\hbar[x^2]=\hbar\). The differential is not a derivation
for the graded product on \(\mathbb R[x]\otimes\Lambda(\varepsilon)\),
where \(|\varepsilon|=-1\) and \(\varepsilon^2=0\): the
Leibniz rule applied to \(x\varepsilon\) would omit the term
\(-\hbar\). Thus a valid reduction of complexes need not be a
reduction of algebras. The specified primitives \(p\varepsilon\)
record which relations have been imposed and how products encounter
them.

Even when integration produces an effective quadratic action, that
action does not determine the normalization. Let \(x\in\mathbb R^n\),
\(y\in\mathbb R^m\), and let
\[
 Q=\begin{pmatrix}A&B\\B^t&C\end{pmatrix}
\]
be real symmetric and positive definite. Here \(C\) is invertible.
Completing the square gives
\begin{equation}\label{v4-prologue:eq:elimination}
 \begin{split}
 &\int_{\mathbb R^m}
   \exp\!\left[-\frac{x^tAx+2x^tBy+y^tCy}{2\hbar}\right]dy\\
 &\qquad=(2\pi\hbar)^{m/2}(\det C)^{-1/2}
   \exp\!\left[-\frac{x^t(A-BC^{-1}B^t)x}{2\hbar}\right].
 \end{split}
\end{equation}
To verify both factors, substitute \(y+C^{-1}B^tx\) for \(y\)
and diagonalize the positive matrix \(C\). Each resulting
one-dimensional integral contributes \(\sqrt{2\pi\hbar/c_i}\),
where \(c_i\) is an eigenvalue of \(C\).
Already with \(B=0\), varying \(C\) leaves the effective action
\(x^tAx/2\) unchanged and changes its amplitude. A comparison of
prepared states must therefore preserve both the quadratic response
and the determinant factor.

\section*{A cubic relation in a strict differential graded algebra}
\label{v4-prologue:cubic}

Let \(k\) be a field of characteristic zero and
\(A=k[x]/(x^3)\), with augmentation \(x\mapsto0\).
A free resolution of the augmentation module \(k\), in nonpositive
cohomological degrees, is
\[
 \cdots\xrightarrow{x}A\xrightarrow{x^2}A
 \xrightarrow{x}A\longrightarrow k\longrightarrow0.
\]
Write \(P^{-n}=Ae_n\) for \(n\ge0\), and
\[
 d_Pe_n=\begin{cases}
 xe_{n-1},&n\text{ odd},\\
 x^2e_{n-1},&n>0\text{ even}.
 \end{cases}
\]
Exactness follows from
\(\ker(x:A\to A)=(x^2)\) and
\(\ker(x^2:A\to A)=(x)\).

The graded endomorphisms of \(P\) form an associative differential
graded algebra \(E\), with composition as product and differential
\[
 D(T)=d_PT-(-1)^{|T|}Td_P.
\]
Define maps \(U,H\) of degree one and \(V\) of degree two by
\[
 \begin{array}{c|cc}
 &n\text{ odd}&n>0\text{ even}\\ \hline
 Ue_n&e_{n-1}&-xe_{n-1}\\
 He_n&0&-e_{n-1}
 \end{array}
 \qquad
 Ve_n=e_{n-2}\quad(n\ge2).
\]
Each map is zero when the indicated output degree is absent.
The two parities give the identities
\begin{equation}\label{v4-prologue:eq:cubic-primitives}
 DU=DV=0,\qquad U^2=-xV=DH,\qquad UH+HU=-V.
\end{equation}
For example, on \(e_2\), both \(U^2\) and \(DH\) give
\(-xe_0\), while \((UH+HU)e_2=-e_0\). The same computations
repeat every two degrees.

Applying \(\operatorname{Hom}_A(-,k)\) to \(P\) gives one copy of
\(k\) in each nonnegative degree and zero differential. Hence
\(\operatorname{Ext}^n_A(k,k)\cong k\) for every \(n\ge0\).
The augmentation \(P\to k\) identifies \(H^*(E)\) with these
groups. Here one uses the elementary property that a bounded-above
complex of projective modules maps every acyclic complex to an
acyclic Hom complex: a cocycle is made null-homotopic successively
from the highest nonzero degree down, lifting at each step through
the surjection onto cycles. This applies to the cone of \(P\to k\).

Write \(u=[U]\) and \(v=[V]\). The maps \(V^r\) and \(UV^r\)
map to nonzero elements of
\(\operatorname{Hom}_A(P^{-2r},k)\) and
\(\operatorname{Hom}_A(P^{-(2r+1)},k)\), respectively.
Also \(UV=VU\), since their coefficients have the same parity.
Together with \(U^2=DH\), this proves
\[
 H^*(E)=\Lambda_k(u)\otimes_k k[v],\qquad |u|=1,\quad |v|=2.
\]

There is more information in \(E\) than in this binary product.
For a degree-one cocycle \(U\) with \(U^2=DH\), the cohomology
class of \(UH+HU\) is its triple Massey product, with this sign
convention. It is closed because
\[
 D(UH+HU)=-U(DH)+(DH)U=0.
\]
Changing either primitive changes its class by an element of
\(uH^1(E)+H^1(E)u\). This indeterminacy is zero here, because
\(H^1(E)=ku\) and \(u^2=0\). Thus
\begin{equation}\label{v4-prologue:eq:massey}
 \langle u,u,u\rangle=\{-v\}\ne\{0\}.
\end{equation}
In the graded algebra \(H^*(E)\) with zero differential, the same
triple product is \(\{0\}\), by choosing both primitives zero.
Maps of differential graded algebras preserve the displayed
construction and its indeterminacy. A quasi-isomorphism therefore
preserves this single-valued product under its cohomology
identification. Consequently \(E\) cannot be joined to
\(H^*(E)\) with zero differential by a chain of quasi-isomorphisms.
Its multiplication was strictly associative from the start.
Strict associativity has retained the cubic information in a
differential and its primitives; it has not removed it.

\section*{An intermediate state in an ordered product}
\label{v4-prologue:matrix}

Let \(e_1,e_2\) be the standard basis of \(k^2\), and let
\(E_{ij}\) send \(e_j\) to \(e_i\) and the other basis vector
to zero. Put \(a=E_{12}\), \(b=E_{21}\). Then
\begin{equation}\label{v4-prologue:eq:ordered-matrix}
 aba=a,\qquad aab=0.
\end{equation}
The first word carries \(e_2\) to \(e_1\), then to \(e_2\),
then to \(e_1\), reading actions from right to left. The reordered
word kills every vector. The order of the intermediate states is
part of the operation.

The normalized trace \(\tau(M)=\operatorname{Tr}(M)/2\) gives
\(\tau(ab)=\tau(ba)=1/2\), and both words in
\eqref{v4-prologue:eq:ordered-matrix} have trace zero. A later
input distinguishes them:
\[
 \tau(abab)=\tfrac12,\qquad \tau(aabb)=0.
\]
Even replacing a two-factor product by the scalar operator with the
same trace changes subsequent composition. The linear projection
\(p(M)=\tau(M)I\) satisfies
\[
 p(ab)a=\tfrac12a\ne aba.
\]
Thus tracing a collision before composing it with the remaining
inputs can discard a coefficient that those inputs still use.

\section*{Positive counting data without the required absolute values}
\label{v4-prologue:synthetic}

Consider the integral matrix
\[
 M=\begin{pmatrix}2&1\\1&3\end{pmatrix},
 \qquad N_n=5^n+1-\operatorname{Tr}(M^n),\qquad n\ge1.
\]
For a positive integer \(d\), let \(\mu(d)\) be the M\"obius
function: it is zero if a prime square divides \(d\), and is
\((-1)^r\) if \(d\) is the product of \(r\) distinct primes.
Define
\begin{equation}\label{v4-prologue:eq:point-counts}
 a_d=\frac1d\sum_{e\mid d}\mu(d/e)N_e.
\end{equation}
Every \(a_d\) is a positive integer. For integrality, regard
\(M\) as the adjacency matrix of a finite directed multigraph.
The number \(\operatorname{Tr}(M^n)\) counts closed walks of
length \(n\) with a marked starting point. Grouping them by their
least period gives
\[
 \operatorname{Tr}(M^n)=\sum_{d\mid n}d\,b_d(M),
\]
where \(b_d(M)\) is the number of primitive cyclic walks of
length \(d\). The same statement holds for the one-vertex
matrices \((5)\) and \((1)\). M\"obius inversion therefore
expresses \(a_d\) as the difference of three integers.

To prove positivity, write
\(\lambda=(5+\sqrt5)/2\) and \(\lambda'=(5-\sqrt5)/2\).
Both are positive, their sum is \(5\), and
\(\operatorname{Tr}(M^n)=\lambda^n+(\lambda')^n\).
Consequently \(0<N_n\le5^n+1\) and
\(N_n\ge5^n-2\lambda^n\).
Every proper divisor of \(d\) is at most \(d/2\), so
\[
 \sum_{\substack{e\mid d\\e<d}}N_e
 \le2\sum_{e=1}^{\lfloor d/2\rfloor}5^e
 <5\cdot5^{d/2}.
\]
Since \(\lambda/5<3/4\), equation
\eqref{v4-prologue:eq:point-counts} gives, for \(d\ge4\),
\[
 da_d>
 5^d\left(1-2(3/4)^d-5^{1-d/2}\right)>0.
\]
The bracket increases with \(d\), and its value at \(d=4\)
is \(107/640\). Direct calculation gives
\[
 a_1=1,\quad a_2=5,\quad a_3=25,
 \quad a_4=110,\quad a_5=500.
\]

Take \(a_d\) formal closed points of degree \(d\). Their zeta
series, defined by a formal Euler product, is
\begin{equation}\label{v4-prologue:eq:synthetic-zeta}
 \begin{split}
 Z(t)&=\prod_{d\ge1}(1-t^d)^{-a_d}
      =\exp\!\left(\sum_{n\ge1}\frac{N_n}{n}t^n\right)\\
     &=\frac{1-5t+5t^2}{(1-t)(1-5t)}.
 \end{split}
\end{equation}
Each coefficient uses only finitely many factors. The middle
identity follows from \(N_n=\sum_{d\mid n}d a_d\). The last
follows by diagonalizing \(M\), or by the formal identity
\(-\log\det(I-tM)=\sum_{n\ge1}\operatorname{Tr}(M^n)t^n/n\).

These counting data admit positive forms of the usual Hall type.
At a point of degree \(d\), take finite-dimensional nilpotent
\(\mathbb F_{5^d}[u]\)-modules. An object \(L\) is a finite
direct sum of such local modules over the formal points. Each
automorphism group is finite. On the complex vector space with
basis their isomorphism classes, the form
\[
 \left\langle\sum_L z_L[L],\sum_L w_L[L]\right\rangle
 =\sum_L\frac{\overline{z_L}w_L}{|\operatorname{Aut}(L)|}
\]
is positive definite on finite sums. Extensions and subobjects are
finite counting problems at each point; no algebraic curve has
been used to supply them.

There is also a reciprocal polynomial and an alternating form:
\[
 P(t)=1-5t+5t^2=5t^2P(1/(5t)),\qquad
 M^tJM=5J,\qquad
 J=\begin{pmatrix}0&1\\-1&0\end{pmatrix}.
\]
Nevertheless the reciprocal roots \(\lambda,\lambda'\) of
\(P\) do not have absolute value \(\sqrt5\).
Positive extension-counting forms and these two identities thus
do not imply the weight-one absolute-value condition. They have
not supplied a geometric Frobenius realization with that property.

The arithmetic comparisons that follow must retain the data exposed
by these calculations. A relation has a specified primitive; a
product can have a nonzero higher operation; a collision can leave
an intermediate state; and a positive auxiliary form need not be
the arithmetic functional to be realized. The comparison maps,
their coefficient domains, and the operations they preserve are
therefore part of the mathematical object studied in this volume.
